\section{Loop 2: Active Learning}

\label{sec:loop2}

\subsection{Motivation}

Users of coding agents routinely correct agent behavior through natural
language. ``No, use tabs not spaces.'' ``Actually, don't use async/await here.''
``Stop importing from lodash.'' These corrections are high-quality preference
signals: they are direct, unambiguous, and immediately actionable. The challenge
is detecting them.

Crucially, these corrections occur within the conversational flow and are not
labeled as corrections by the user. An LLM asked to identify corrections in a
transcript will find them, but with high false-positive rate (flagging
instructions as corrections) and meaningful false-negative rate (missing
implicit corrections that come as rewrites rather than explicit statements).
We propose a structural detection approach that is faster and more reliable.

\subsection{Structural Signal Detection}

\begin{definition}[Correction Signal]
\label{def:correction}
A correction signal $\sigma$ is detected when any of the following structural
conditions hold in the most recent $W=5$ turns:
\begin{itemize}[leftmargin=1.5em, nosep]
    \item \textbf{Edit-distance trigger:} The user's latest message $m_t$ has
          normalized Levenshtein distance $d_L(m_t, m_{t-2}) / |m_{t-2}| < 0.3$,
          where $m_{t-2}$ is the agent's preceding response (indicating direct
          text editing of the agent's output).
    \item \textbf{Negation trigger:} $m_t$ begins with or contains one of
          the negation phrases in $\mathcal{N}$ (e.g., ``no,'' ``actually,''
          ``don't,'' ``stop,'' ``instead'').
    \item \textbf{Tool override trigger:} The user explicitly invokes a tool
          with parameters that directly contradict the agent's last tool call
          on the same target.
    \item \textbf{Rewrite trigger:} The user provides a full rewrite of
          agent-generated code that is $>50\%$ different by line count.
\end{itemize}
\end{definition}

\subsection{LearnedFact Schema}

When a correction signal $\sigma$ is detected, the active learning loop
extracts a \textsc{LearnedFact} record and appends it to the facts database.

\begin{definition}[LearnedFact Record]
\label{def:learned_fact}
A \textsc{LearnedFact} record $\phi$ is a tuple:
\begin{equation}
\phi = \bigl( \text{id},\; \text{content},\; \text{confidence},\;
              \text{scope},\; \text{source},\; \text{ts} \bigr),
\end{equation}
where:
\begin{itemize}[leftmargin=1.5em, nosep]
    \item $\text{content} \in \Sigma^*$: Natural language statement of the fact.
    \item $\text{confidence} \in [0, 1]$: Prior confidence in the fact.
    \item $\text{scope} \in \{\text{global}, \text{project}, \text{session}\}$: Applicability scope.
    \item $\text{source} \in \{\text{user-stated}, \text{inferred}, \text{corrected}\}$: Provenance.
\end{itemize}
\end{definition}

A critical design decision is that user-stated facts---facts extracted from
a negation trigger where the user explicitly states a preference---receive
$\text{confidence} = 1.0$ and $\text{scope} = \text{project}$ by default.
This reflects the epistemic status of direct user instruction: it is not a
hypothesis to be tested but a ground truth to be respected.

Inferred facts (extracted from edit-distance or rewrite triggers) receive
$\text{confidence} = 0.7$ initially, which is updated by subsequent
confirmations or contradictions. A conflicting user-stated fact always
supersedes an inferred fact, regardless of confidence.

\subsection{Why Hooks Beat LLM Judgment}

\begin{proposition}[Structural Detection Superiority]
\label{prop:structural_detection}
For correction detection in coding agent sessions, structural signal detection
achieves precision $P_s \geq P_{\text{LLM}}$ and recall $R_s \geq R_{\text{LLM}}$
with time complexity $O(1)$ per turn compared to $O(n)$ for LLM-based detection
(where $n$ is context length).
\end{proposition}

The argument is empirical: the structural triggers in Definition~\ref{def:correction}
correspond directly to the linguistic behaviors users exhibit when correcting
agents. The negation trigger alone captures $>80\%$ of explicit corrections.
The remaining triggers capture implicit corrections that LLMs systematically
miss because they require comparing tool call arguments to prior tool call
arguments---a structured comparison that LLMs perform unreliably but code
performs exactly.

% ============================================================================
% 7. LOOP 3: SESSION REFLECTION
% ============================================================================
